\section{Conclusion}
\label{sec:conclusion}

A pre-registered sweep over $\tau \in \{0.02, 0.03, 0.05, 0.075,
0.10\}$ across 11{,}250 result rows produces a sharp negative
feasibility result: \textbf{no $\tau$ in the studied grid
simultaneously satisfies Paper~10's two pre-registered tolerances.}
The K=50 adaptive--lag1 cost is between $-4.1$ and $-5.3$~pp at
every $\tau$, always exceeding the $-2.0$~pp ceiling; the K=200
worst-case adaptive--fixed delta is uniformly $0.0$. The simple
magnitude detector is fundamentally insufficient for the joint
feasible region.

The H1 finding is stronger than Paper~10 reported: the K=200
hazard is uniformly avoided across an order-of-magnitude $\tau$
range, not merely at the specific $\tau = 0.05$ Paper~10 used.

The H2 and H4 hypothesis directions were pre-registered with the
wrong sign; the data report Spearman correlations of opposite sign
to the protocol's prediction ($\rho_{H4} = -0.97$, $\rho_{H2} =
+0.87$). The honest interpretation is that ``larger $\tau$ fires
less, not more,'' and that the K=50 cost shrinks with $\tau$, not
grows. The pre-registration discipline correctly surfaces this
mathematical sign error as a falsified prediction rather than
allowing post-hoc correction.

\textbf{The deployable conclusion:} simple adaptive switching adds
operational complexity without reaching Paper~10's joint feasibility
region across the studied $\tau$ range. For deployments with stable
capacity, the static rule (lag-1 if $K \leq 100$, fixed if
$K \geq 200$) is sufficient and adaptive procedures should be
preferred only when a richer detector --- capacity-aware
threshold, CUSUM, or Bayesian online change-point --- has been
evaluated under pre-registration.

\textbf{Reproducibility.} Runner, analysis, pre-registration protocol,
and frozen 11{,}250-row results CSV are available at
\url{https://github.com/Harshavardhanmalla-Labs/research-papers/tree/main/paper11}.

\textbf{Future work.} (i) Capacity-aware-threshold detector with
$\tau_K$ pre-registered per capacity; (ii) CUSUM~\cite{page1954cusum}
detector with pre-registered accumulator parameters; (iii) Bayesian
online change-point detector; (iv) sensitivity sweep over the
calibration grid coarseness; (v) real fleet validation with
longitudinal exploitation outcome data.
